\chapter{Fundamentare Teoretică}
\label{cap:fundamentare_teoretica}

Acest capitol detaliază fundamentele teoretice și algoritmice care stau la baza sistemului StyleGenAI. Sunt prezentate arhitectura software, modelul de date utilizat în MVP, pipeline-ul de procesare a imaginilor și logica motorului de recomandare. Accentul este pus pe coerența dintre descrierea teoretică și implementarea efectivă a aplicației.

\section{Arhitectura Sistemului și Modelul de Date}
\label{sec:arhitectura_cap4}

Așa cum a fost introdus în capitolele anterioare, sistemul utilizează o arhitectură Client-Server, decuplând interfața de utilizator (client) de logica de procesare (server).

\begin{itemize}
    \item \textbf{Client (Frontend)}: O aplicație mobilă cross-platform dezvoltată în React Native și Expo. Responsabilitățile sale includ randarea interfeței grafice, gestionarea interacțiunilor cu utilizatorul (încărcare imagini, selecție categorii) și comunicarea cu backend-ul prin API REST.
    \item \textbf{Server (Backend)}: O aplicație API RESTful construită cu Python și framework-ul Flask. Aceasta orchestrează întreaga logică de business: autentificarea utilizatorilor, stocarea și procesarea imaginilor, generarea și evaluarea ținutelor, și persistența datelor.
\end{itemize}

Datele sunt interschimbate în format JSON, iar autentificarea și autorizarea cererilor se realizează prin JSON Web Tokens (JWT). În mediul de producție, comunicația prin HTTPS este tratată la nivel de infrastructură (de exemplu, reverse proxy), nu impusă rigid în logica internă a aplicației.

\subsection{Modelul de Date}
Persistența datelor este asigurată, în configurația de MVP, de o bază de date relațională SQLite, aleasă pentru simplitatea configurării și portabilitate. Aplicația include și suport pentru SQL Server, însă modelul logic rămâne similar. În varianta curentă, schema se concentrează pe următoarele entități:

\begin{itemize}
    \item \texttt{users}: Stochează datele de autentificare și rolul utilizatorului (email, hash de parolă, rol, timestamp-uri).
    \item \texttt{outfits}: Stochează istoricul ținutelor/sugestiilor salvate pentru utilizator (\texttt{image\_url}, \texttt{style\_data}, \texttt{liked}, \texttt{created\_at}).
    \item \texttt{ai\_metrics}: Memorează metadate operaționale despre generări (stil, sezon, gen, timp de procesare, succes/eroare), utile pentru monitorizare.
\end{itemize}

În implementarea actuală nu există o tabelă separată \texttt{wardrobe\_items}. Datele relevante pentru piese și combinații sunt gestionate în fluxul de generare și în câmpurile serializate ale istoricului.

\section{Pipeline-ul de Procesare a Imaginilor}
\label{sec:image_pipeline}

Un pilon central al aplicației este capacitatea de a transforma o imagine brută a unui articol vestimentar într-un set de atribute structurate. Acest proces se desfășoară în două etape majore, executate pe backend.

\begin{figure}[h]
    \includegraphics[width=\textwidth]{imagini/fluxul imaginii.png}
    \caption{Fluxul de procesare a imaginii: de la imaginea originală la extragerea culorii dominante.}
    \label{fig:image_pipeline}
\end{figure}

\subsection{Segmentarea Obiectului de Interes}
Primul pas constă în izolarea articolului vestimentar de fundalul imaginii. O fotografie realizată de utilizator poate conține elemente irelevante (ex: umeraș, perete, alte obiecte) care pot influența negativ analiza culorii. Pentru a reduce acest efect, backend-ul încearcă aplicarea unei etape de eliminare a fundalului prin biblioteca \texttt{rembg}.

\texttt{rembg} folosește modele pre-antrenate de detecție a obiectului salient, inclusiv familia U²-Net \cite{qin2020u2net}. În aplicația curentă, această etapă este \textbf{opțională}: dacă procesarea reușește, se obține o versiune cu fundal transparent; dacă eșuează, fluxul continuă cu imaginea originală. Astfel, pipeline-ul rămâne robust și nu blochează utilizatorul în caz de erori locale de procesare.

\subsection{Extragerea Culorii Dominante cu K-Means}
În implementarea de MVP, extragerea culorii dominante este realizată printr-o metodă euristică simplă, eficientă și ușor de explicat. După deschiderea imaginii în format RGBA, se păstrează doar pixelii netransparenți (alpha $>$ 0), iar culoarea dominantă este aproximată prin media aritmetică a canalelor RGB pentru acești pixeli.

Formal, dacă $P=\{(r_i,g_i,b_i)\}_{i=1}^{n}$ este mulțimea pixelilor relevanți, culoarea estimată este:

\begin{equation}
\bar{c} = \left(\frac{1}{n}\sum_{i=1}^{n} r_i,\, \frac{1}{n}\sum_{i=1}^{n} g_i,\, \frac{1}{n}\sum_{i=1}^{n} b_i\right)
\end{equation}

Rezultatul numeric este completat prin maparea la cea mai apropiată etichetă de culoare dintr-un set predefinit (ex: \textit{negru}, \textit{alb}, \textit{gri}, \textit{bej}, \textit{bleumarin}). Această soluție favorizează simplitatea, predictibilitatea și costul computațional redus, fiind adecvată pentru un prototip funcțional.

\section{Fundamentele Teoriei Culorilor}
\label{sec:color_theory}

Motorul de recomandare se bazează pe principii din teoria culorilor pentru a evalua compatibilitatea dintre articole. Pentru a înțelege deciziile algoritmice, este necesară o scurtă introducere în conceptele de bază.

\subsection{Spații de Culoare: RGB vs. HSV}
\begin{itemize}
    \item \textbf{RGB (Red, Green, Blue)}: Este un model de culoare aditiv, unde culorile sunt create prin combinarea luminii roșii, verzi și albastre. Este modelul nativ pentru ecranele digitale. Totuși, este un spațiu ne-perceptual, ceea ce înseamnă că distanța euclidiană între două culori în RGB nu corespunde întotdeauna cu similaritatea percepută de ochiul uman.
    \item \textbf{HSV (Hue, Saturation, Value)}: Este un model de culoare cilindric, mai intuitiv pentru percepția umană.
        \begin{itemize}
            \item \textbf{Hue (Nuanța)}: Reprezintă tipul pur al culorii (roșu, verde, albastru) și este măsurat ca un unghi (0-360 de grade) pe roata culorilor. Este cea mai importantă componentă pentru a determina relațiile de armonie.
            \item \textbf{Saturation (Saturația)}: Reprezintă intensitatea sau puritatea culorii (0-100\%). O saturație de 0\% corespunde unei nuanțe de gri.
            \item \textbf{Value (Valoarea/Luminozitatea)}: Reprezintă luminozitatea culorii (0-100\%). O valoare de 0\% este negru.
        \end{itemize}
\end{itemize}

HSV este adesea preferat în sisteme orientate strict pe armonii cromatice perceptuale. În schimb, implementarea actuală StyleGenAI utilizează o euristică în RGB pentru a păstra algoritmul simplu și transparent în stadiul de MVP.

\begin{figure}[h]
    \centering
    \includegraphics[width=\textwidth]{imagini/comparatia vizuala.png}
    \caption{Comparație vizuală între modelul cartezian RGB (stânga) și modelul cilindric HSV (dreapta).}
    \label{fig:rgb_hsv}
\end{figure}

\subsection{Armonii Cromatice}
Pe baza roții culorilor, se definesc mai multe scheme de armonie, dintre care cele mai relevante pentru proiect sunt:
\begin{itemize}
    \item \textbf{Culori Analoage}: Culori care sunt adiacente pe roata culorilor (ex: roșu, roșu-oranj, oranj). Creează un efect vizual plăcut și calm.
    \item \textbf{Culori Complementare}: Culori care sunt opuse pe roata culorilor (ex: roșu și verde). Creează un contrast puternic și vibrant.
    \item \textbf{Culori Monocromatice}: Variații de saturație și luminozitate ale aceleiași nuanțe.
    \item \textbf{Culori Neutre}: Culori precum alb, negru, gri și bej, care se pot combina ușor cu majoritatea paletelor.
\end{itemize}

\section{Motorul de Recomandare (Styling Engine)}
\label{sec:recomandare_cap4}

Sistemul de recomandare este unul de tip \textbf{content-based}, deoarece ia decizii exclusiv pe baza atributelor intrinseci ale articolelor (culoare, categorie), fără a se baza pe interacțiunile altor utilizatori.

\subsection{Generarea Spațiului de Căutare}
Pentru un utilizator cu seturile de articole $T$ (Topuri), $B$ (Pantaloni) și $S$ (Încălțăminte), spațiul total al ținutelor posibile este produsul cartezian $T \times B \times S$. Algoritmul implementează o abordare de tip forță brută, generând explicit fiecare combinație cu \texttt{itertools.product}. Complexitatea rămâne $O(|T| \cdot |B| \cdot |S|)$.

Înainte de generare, fluxul aplică filtre de stil și verifică o condiție minimă de inventar (minimum 5 articole pe fiecare categorie principală), pentru a evita recomandări triviale sau instabile.

\subsection{Funcția de Scor pentru Compatibilitate Cromatică}
Fiecare ținută generată este evaluată printr-o funcție de scor euristică aplicată culorilor în format hex/RGB. În implementarea actuală, scorul agregă mai multe componente: armonie pe baza distanțelor RGB, balans de contrast, mix neutru-accent, bonus de diversitate, penalizare pentru repetiție, plus ajustări ușoare după stil și sezon.

\begin{equation}
S \approx S_{harmony} + S_{contrast} + S_{accent} + S_{diversity} - P_{repetition} + M_{style} + M_{season}
\label{eq:compatibility_score_rgb}
\end{equation}

Măsurarea similarității între culori folosește distanța euclidiană normalizată în RGB:

\begin{equation}
d(c_a, c_b)=\frac{\sqrt{(R_a-R_b)^2+(G_a-G_b)^2+(B_a-B_b)^2}}{\sqrt{255^2+255^2+255^2}}
\end{equation}

\begin{itemize}
    \item \textbf{Armonie ($S_{harmony}$)}: Componentă principală derivată din distanțele dintre perechi de culori în RGB, favorizând un echilibru cromatic moderat.
    
    \item \textbf{Contrast și mix neutru-accent}: Se acordă punctaj în funcție de numărul perechilor suficient de distincte și de prezența culorilor neutre în combinație.
    
    \item \textbf{Diversity/penalty}: Se favorizează combinațiile variate, dar se penalizează repetiția excesivă a aceleiași culori.
    
    \item \textbf{Modificatori contextuali}: Stilul cerut și sezonul pot adăuga ajustări mici scorului final.
    
    \item \textbf{Semnal de trend}: După calculul scorului de bază, generatorul aplică un bonus fix dacă cel puțin o culoare a ținutei se regăsește în lista de culori populare definită în fișierul de trenduri.
\end{itemize}

După evaluarea tuturor combinațiilor, ținuta cu scorul maxim este selectată și returnată utilizatorului împreună cu o analiză a componentelor scorului. În plus, pentru garderobe mai mari, sistemul utilizează un mecanism de cache pentru a evita recalculările redundante.
